\section{Fermat 与 Euler 定理}
\label{sec:04-Discrete-Mathematics/04-Number-Theory/03}

\subsection{一个计算困境}

算 \(3^{100}\pmod 7\)。你若试图先算出 \(3^{100}\)——它是一个 48 位的十进制数——再除以 7 取余数，这条路既笨又不必要。困难在于，我们需要的只是余数，但指数把中间量撑到了天文数字的规模。

上两节已经建立了模运算的基本规则：加、减、乘在模 \(n\) 下可以放心操作，上一节还得到了消去律——当 \(\gcd(a,n)=1\) 时，\(ar\equiv as\) 意味着 \(r\equiv s\)。这些工具对付一次乘法绰绰有余，可一百次乘法呢？快速幂能把 100 次压缩到大约 \(\log_2 100\approx 7\) 轮平方与乘法，效率够了，但我们想要的不是更快地算出 \(3^{100}\bmod 7\)，而是理解指数幂在模的世界里到底遵循什么规律——有没有一个隐藏的周期，让巨大的指数塌缩成很小的等价物？

\subsection{一个小发现：指数确实会循环}

动手算前几个幂：\(3^1\equiv3\)，\(3^2\equiv2\)，\(3^3\equiv6\)，\(3^4\equiv4\)，\(3^5\equiv5\)，\(3^6\equiv1\pmod 7\)。到第 6 次方回归到 1。一旦某次方等于 1，后面的幂就循环了：\(3^7\equiv3\)，\(3^8\equiv2\)，依此类推。所以指数模 6 就够了：\(3^{100}=3^{6\cdot16+4}\equiv 3^4=81\equiv 4\pmod 7\)。

回过头看，循环长度是 6，正好是 \(\varphi(7)=7-1\)。这只是一个巧合，还是触及了某种结构？

试探另一个模数：模 10。与 10 互素的余数是 \(\{1,3,7,9\}\)，共 \(\varphi(10)=4\) 个。取 \(a=3\)：\(3^1=3\)，\(3^2=9\)，\(3^3=27\equiv7\)，\(3^4=81\equiv1\pmod{10}\)。又是 4——又等于 \(\varphi(10)\)。再试模 15：\(\varphi(15)=8\)，取 \(a=2\)：\(2^1=2\)，\(2^2=4\)，\(2^3=8\)，\(2^4=16\equiv1\pmod{15}\)——周期只有 4，比 \(\varphi(15)=8\) 短。所以 \(\varphi(n)\) 不是一个精确周期，但它像是一个``上界''——指数在到达 \(\varphi(n)\) 之前或之时，一定会回到 1。

\subsection{把直觉变成证明：乘法置换互素余数}

设 \(n\) 固定，列出所有与 \(n\) 互素的余数类：
\[
r_1,r_2,\ldots,r_{\varphi(n)}.
\]
这 \(\varphi(n)\) 个数每个都可逆（存在模 \(n\) 下的乘法逆元）。现在取一个与 \(n\) 互素的 \(a\)，把上面每个数都乘上 \(a\)：
\[
ar_1,ar_2,\ldots,ar_{\varphi(n)}.
\]
乘出来的这 \(\varphi(n)\) 个数，仍然是全部互素余数类——只是顺序变了。

理由分两步。第一，每个 \(ar_i\) 仍然与 \(n\) 互素，因为 \(a\) 和 \(r_i\) 都与 \(n\) 互素。第二，它们两两不同：若 \(ar_i\equiv ar_j\pmod n\)，由消去律（\(\gcd(a,n)=1\)）立刻得到 \(r_i\equiv r_j\)，即 \(i=j\)。一个有限集合上的单射必然也是满射——所以 \(\{ar_i\}\) 就是 \(\{r_i\}\) 的一个重新排列。

现在把两个集合的元素全部乘起来。乘积相等：
\[
(ar_1)(ar_2)\cdots(ar_{\varphi(n)})
\equiv r_1r_2\cdots r_{\varphi(n)}\pmod n.
\]
左边提出 \(\varphi(n)\) 个 \(a\)：
\[
a^{\varphi(n)}\cdot r_1r_2\cdots r_{\varphi(n)}
\equiv r_1r_2\cdots r_{\varphi(n)}\pmod n.
\]
令 \(P=r_1r_2\cdots r_{\varphi(n)}\)。\(P\) 中每个因子都与 \(n\) 互素，所以 \(P\) 也与 \(n\) 互素，在模 \(n\) 下可逆。两边消去 \(P\)，得到 Euler 定理：
\[
a^{\varphi(n)}\equiv1\pmod n,\qquad\gcd(a,n)=1.
\]

当 \(p\) 为素数时，\(\varphi(p)=p-1\)，此时定理退化为 Fermat 小定理：
\[
a^{p-1}\equiv1\pmod p,\qquad p\nmid a.
\]
两边同乘 \(a\) 得到等价的 \(a^p\equiv a\pmod p\)，这个形式对 \(p\mid a\) 也成立（两边同为 \(0\)）。

\subsection{回到数字：这个定理到底怎么用}

回到最初的 \(3^{100}\pmod 7\)。\(\varphi(7)=6\)，Euler 定理保证 \(3^6\equiv1\pmod 7\)。用除法把指数拆开：\(100=16\cdot6+4\)，于是
\[
3^{100}=(3^6)^{16}\cdot3^4\equiv1^{16}\cdot3^4=81\equiv4\pmod 7.
\]
这就是前面发现``指数模 6''的道理——Euler 定理告诉我们，底数与模数互素时，指数可以模 \(\varphi(n)\) 约化。你再算 \(3^{1000}\pmod 7\)，也只消做一次 \(1000\bmod 6=4\)，\(3^4=81\equiv4\)——两秒的事。

但约化指数之前必须确认底数与模数互素。以 \(n=4\) 为例，\(\varphi(4)=2\)，取 \(a=2\)：\(\gcd(2,4)=2\neq1\)，直接代公式会给出 \(2^2=4\equiv0\not\equiv1\pmod 4\)。定理的前提不是装饰——不互素时乘法置换的论证从第一步就崩塌了，因为 \(a\) 不能消去。

另一个注意点：Euler 定理给的是周期的上界，不保证最小周期。前面模 15 的例子中，\(2^4\equiv1\) 而 \(\varphi(15)=8\)。实际乘法阶（使 \(a^k\equiv1\pmod n\) 的最小正整数 \(k\)）可以是 \(\varphi(n)\) 的真因子。定理对任意 \(a\) 只保证``到达 \(\varphi(n)\) 时一定回到 1''，至于更早回来，那是额外的运气。

\subsection{一个危险的推论：Fermat 测试不是素数证明}

Fermat 小定理提供了一个素性检测思路：对候选的数 \(n\)，随便挑一个 \(a\) 满足 \(\gcd(a,n)=1\)，检验 \(a^{n-1}\equiv1\pmod n\)。若不成立，\(n\) 一定是合数——这是一条铁证。但反过来呢？若对某个 \(a\) 成立，能不能宣告 \(n\) 是素数？

不能。存在合数 \(n\) 和底数 \(a\) 满足 \(a^{n-1}\equiv1\pmod n\)——这些 \(n\) 对底数 \(a\) 而言是伪素数（pseudoprime）。更麻烦的是 Carmichael 数：对\textbf{所有}与 \(n\) 互素的底数 \(a\) 都通过 Fermat 测试的合数。最小的 Carmichael 数是 561，它是合数（\(561=3\cdot11\cdot17\)），但对任意 \(\gcd(a,561)=1\) 都有 \(a^{560}\equiv1\pmod{561}\)。

Fermat 测试只能提供单边证据：没通过，一定是合数；通过了，只能说是``候选素数''。可靠素性检验需要更强的工具（Miller--Rabin 等），这已超出本单元范围。重点是理解定理自身的边界：它给你一个对所有互素底数都成立的同余式，但同余式的成立反过来不保证模数是素数。

\subsection{RSA 的数学骨架}

RSA 把 Euler 定理变成了加密工具。选大素数 \(p,q\)，令 \(n=pq\)，则 \(\varphi(n)=(p-1)(q-1)\)。选加密指数 \(e\) 与 \(\varphi(n)\) 互素，求解密指数 \(d\) 满足
\[
ed\equiv1\pmod{\varphi(n)}.
\]
即存在整数 \(k\) 使 \(ed=1+k\varphi(n)\)。

加密：\(c\equiv m^e\pmod n\)。解密：\(c^d\equiv m^{ed}\pmod n\)。当 \(\gcd(m,n)=1\) 时，Euler 定理直接保证 \(m^{ed}=m\cdot(m^{\varphi(n)})^k\equiv m\pmod n\)。当 \(m\) 与 \(n\) 不互素（即 \(m\) 是 \(p\) 或 \(q\) 的倍数）时，需要分别在模 \(p\) 和模 \(q\) 下验证再用中国剩余定理拼回——这是下节的内容。安全性来自从公开的 \(n\) 恢复其因子在计算上的困难，而不是模幂运算本身难算。

这里展示的只是数学骨架。真实密码系统还需要安全填充方案、足够的密钥长度、侧信道防护等工程条件。直接用``教科书 RSA''处理真实消息是不安全的——比如相同的 \(m\) 和 \(e\) 在不同接收者之间可能泄漏明文，短消息的加密可能被字典攻击枚举。但骨架本身已经把 Euler 定理从一段漂亮的数学变成了一个能锁门的函数，这是值得体会的。

若希望把``可逆余数类''的结构进一步抽象为群的实例，可继续阅读\hyperref[ch:032]{``群与对称：可逆操作怎样组织不变性''}；本节的置换证明正是有限群中 Lagrange 型周期现象在数论中的具体投影。
